% Source: https://artofproblemsolving.com/community/c3h2790818_2019_mbmt_guts_round_d1630_l1015_montgomery_blair_math_tournament
% Stable identity: aops-topic-2790818-q23
Describe the solid by the heights $h(i,j)$ of its $25$ columns over a $5\times5$ base. The top view forces every $h(i,j)\ge1$; the front and side views force the tallest column in each row and in each column of the base to have height $5$.

The largest volume is $125$ (the full cube). For the smallest, the height-$5$ columns must meet every row and column of the base, so there are at least $5$ of them, and the other columns have height at least $1$: volume at least $5\cdot5+20=45$. This is achieved with the tall columns on a diagonal (heights shown from above):
\begin{asy}
usepackage("amsmath");
size(3cm);
defaultpen(fontsize(10pt));
for(int i=0;i<5;++i) for(int j=0;j<5;++j){draw(shift(i,j)*unitsquare);label((i==j)?"$5$":"$1$",(i+0.5,j+0.5));}
\end{asy}

The difference is $125-45=\boxed{80}$.
